% !TEX TS-program = lualatex

\DocumentMetadata{
  lang = en,
  pdfversion = 2.0,
  pdfstandard = ua-2,
  tagging = on,
  tagging-setup = {
    math/setup = {mathml-SE},
    table/header-rows = 1
  }
}

\documentclass[12pt]{article}

\usepackage{amsmath}
\usepackage{geometry}
\geometry{
  letterpaper,
  textwidth=6.5in,
  textheight=9in,
  top=0.85in,
  bottom=1in,
  left=0.75in,
  right=0.75in,
}
\usepackage{hyperref}

\newtheorem{theorem}{Theorem}
\newtheorem{acknowledgement}[theorem]{Acknowledgement}
\newtheorem{algorithm}[theorem]{Algorithm}
\newtheorem{axiom}[theorem]{Axiom}
\newtheorem{case}[theorem]{Case}
\newtheorem{claim}[theorem]{Claim}
\newtheorem{conclusion}[theorem]{Conclusion}
\newtheorem{condition}[theorem]{Condition}
\newtheorem{conjecture}[theorem]{Conjecture}
\newtheorem{corollary}[theorem]{Corollary}
\newtheorem{criterion}[theorem]{Criterion}
\newtheorem{definition}[theorem]{Definition}
\newtheorem{example}[theorem]{Example}
\newtheorem{exercise}[theorem]{Exercise}
\newtheorem{lemma}[theorem]{Lemma}
\newtheorem{notation}[theorem]{Notation}
\newtheorem{problem}[theorem]{Problem}
\newtheorem{proposition}[theorem]{Proposition}
\newtheorem{remark}[theorem]{Remark}
\newtheorem{solution}[theorem]{Solution}
\newtheorem{summary}[theorem]{Summary}
\newenvironment{proof}[1][Proof]{\textbf{#1.} }{\ \rule{0.5em}{0.5em}}

\tagpdfsetup{
  math/alt/use
}

\hypersetup{
  pdftitle={Assignment 7},
  pdfauthor={Blankenau},
  pdfsubject={Fall, 2025},
  hidelinks,
}

\begin{document}
\begin{center}
{\LARGE Assignment 7}
\end{center}

\noindent Blankenau

\noindent Economics 905, Fall, 2025

\medskip
\noindent \noindent \textit{Instructions.} You may work in groups and may
compare answers prior to submission. Joint submissions with up to 4 names
are allowed and encouraged. All work should be legible and concise. A due
date will be announced in class.

\medskip Consider a discrete time endogenous growth model where a
representative household maximizes%
\begin{equation*}
U=\overset{\infty }{\underset{t=0}{\sum }}\beta ^{t}\frac{C_{t}^{\sigma }}{%
\sigma }
\end{equation*}%
$0<\beta <1\ $where $C_{t}$ is consumption. The population is normalized to
1. Each agent has one unit of time available in each period to allocate
across human capital accumulation and production. The agents can produce
consumption goods according to $Y_{t}=Bh_{t}u_{t}$ where $Y_{t}$ is the
output of the final good$.$ Note, then, that $\left( 1-u_{t}\right) $ is the
time spent in human capital accumulation. The human capital accumulation
technology is given by $h_{t+1}=Ah_{t}\left( 1-u_{t}\right) .\ $where $A>0.$

\begin{enumerate}
\item Set up the agent's optimization problem as a dynamic programming
problem and find the first order condition(s) and envelope condition(s).

\item Find the balanced growth rate of consumption in this model and find $u$
along the balanced growth path.

\item Let the period utility function be $\ln C_{t}.$ Now suppose there is
congestion in schooling so that the more time people spend studying the less
productive is the educational system. This is expressed by $A=\left( 1-\bar{u%
}\right) ^{-\alpha }$ where $0<\alpha <1$ and $\bar{u}$ is the average
amount of time spent in schooling, taken as given by all agents but equal to 
$u$ in equilibrium. Find the balanced growth rate of consumption in this
model and find $u$ along the balanced growth path.

\item With again $A=\left( 1-\bar{u}\right) ^{-\alpha },$ write down and the
social planner's problem and find first order conditions and envelope
conditions. Recall the social planner internalizes the externality and so
sets $A=(1-u_{t})^{-\alpha }$.

\item Show that the social planners's allocation would yield%
\begin{equation*}
\gamma =\left( \frac{\beta \left( 1-\alpha \right) }{1-\beta \alpha }\right)
^{1-\alpha }
\end{equation*}%
\begin{equation*}
u=1-\frac{\beta \left( 1-\alpha \right) }{1-\beta \alpha }.
\end{equation*}
\end{enumerate}
\end{document}
